% GENERATED by scripts/manual_tables.py from the adapter registry. Do not edit:
% rerun the script. Hand-written text: blocks/detection/impulse_events/
\ubblock{detection.impulse_events}%
  {Impulse events}% what the app calls it
  {Impulse events: one event per impulse, running to the next}% registry display name
  {detect}% insert-modal tab
  {Scores}% input
  {SpanSet}% output
  {Adapters/detection\_impulse\_events.py}% code
  {detection/impulse_events}% hand-written folder under blocks/
\ubblockdescription{detection.impulse_events}{Turns an impulse-like score (mostly zero, sharp impulses --- e.g. Deconvolve's jumps) into one event per impulse: from pre\_s before the jump to the next jump or max\_extent\_s, whichever is first, so an event that starts before the last one ended is still its own event. Tune threshold first (drag it on the score); merge\_within\_s joins a fall spread over a few samples. Each event's onset and trough are marked, written to `detections`.}
\ubblockflags{detection.impulse_events}{%
}
\ubblocksideinputs{detection.impulse_events}{%
}
\ubblockparams{detection.impulse_events}{%
  \ubparam{threshold}{4.0}{number, at least 0}{The cut, in the score's units (for Deconvolve: the jump size in noise SDs); an impulse is a sample above it}%
  \ubparam{merge\_\allowbreak{}within\_\allowbreak{}s}{2.0}{number, at least 0}{Impulses closer than this many seconds are one event (a fall spread over a few samples)}%
  \ubparam{pre\_\allowbreak{}s}{10.0}{number, at least 0}{Seconds of context kept before each event's onset}%
  \ubparam{max\_\allowbreak{}extent\_\allowbreak{}s}{600.0}{number, at least 1}{The longest an event may run after its onset when no next event ends it first}%
}
